% Part: computability
% Chapter: recursive-functions
% Section: non-pr-functions

\documentclass[../../../include/open-logic-section]{subfiles}

\begin{document}

\olfileid{cmp}{rec}{npr}
\olsection{非原始递归函数}

原始递归函数并未穷尽直观上可计算的函数。直观上应该清楚，我们可以列出所有一元原始递归函数 $f_0$, $f_1$, $f_2$, \dots，使得我们可以能行地计算 $f_x$ 在输入 $y$ 上的值；换句话说，由
\[
g(x,y) = f_x(y)
\]
定义的函数 $g(x,y)$ 是可计算的。进而函数
\begin{eqnarray*}
h(x) & = & g(x,x) + 1 \\
& = & f_x(x) +1.
\end{eqnarray*}
也是可计算的。对于每个原始递归函数 $f_i$，$h$ 与 $f_i$ 在 $i$ 处的值不同。所以 $h$ 是可计算的，但不是原始递归的；对 $g$ 亦然。这是 康托尔对角线论证的一个“能行”版本。

我们可以给出更为明确的非原始递归可计算函数的例子。例如，令记法 $g^n(x)$ 表示将 $g$ 迭代 $n$ 次作用于 $x$，即$g(g(\dots g(x)))$；并通过
\begin{eqnarray*}
g_0(x) & = & x+1 \\
g_{n + 1}(x) & = & g_n^x(x)
\end{eqnarray*}
定义一个函数序列 $g_0, g_1, \dots$。可以验证每个函数 $g_n$ 都是原始递归的。序列中相继的每个函数都比前一个增长得快得多；$g_1(x)$ 等于 $2x$，$g_2(x)$ 等于 $2^x \cdot x$，而 $g_3(x)$ 的增长大致相当于 $x$ 个 $2$ 组成的指数塔。Ackermann--P\'eter 函数本质上就是函数$G(x) = g_x(x)$，并且可以证明它的增长速度超过任何原始递归函数。

让我们回到枚举原始递归函数的问题。请记住，我们已经为每个原始递归函数分配了符号记法；因此，枚举这些记法就足够了。我们可以如下递归地为每个记法 $F$ 分配一个自然数 $\#(F)$：
\begin{eqnarray*}
\#(0) & = & \langle 0 \rangle \\
\#(S) & = & \langle 1 \rangle \\
\#(\Proj{n}{i}) & = & \langle 2, n, i \rangle \\
\#(\fn{Comp}_{k,l}[H,G_0,\dots,G_{k-1}]) & = & \langle
3,k,l,\#(H),\#(G_0),\dots,\#(G_{k-1}) \rangle \\
\#(\fn{Rec}_l[G,H]) & = & \langle 4, l, \#(G), \#(H) \rangle
\end{eqnarray*}
这里我们利用了上一节的编码，使得每个数字序列都可以被视为一个自然数。由此可知，每个记法都被赋予了一个自然数编码。当然，有些序列（因此有些数）并不对应任何记法；但我们可以定义 $f_i$ 为：如果 $i$ 是某个记法的编码，则 $f_i$ 就是以该记法表示的一元原始递归函数，否则 $f_i$ 就是常值 $0$ 函数。最终结果是，我们有了一个显式枚举一元原始递归函数的方法。

（事实上，有些函数，比如常值 $0$ 函数，会在列表中出现不止一次。这不仅是我们编码方式造成的，也源于常值 $0$ 函数本身就有不止一种记法这一事实。稍后我们会看到，无法可计算地避免这些重复；例如，不存在可计算函数来判定一个给定的记法是否表示常值 $0$ 函数。）

现在我们可以令函数 $g(x,y)$ 由 $f_x(y)$ 给出，这里 $f_x$ 指代我们刚刚描述的枚举中的函数。我们如何知道 $g(x,y)$ 是可计算的呢？直观上很清楚：要计算 $g(x,y)$，首先“解包” $x$，看看它是否是一个一元函数的记法。如果是，就计算该函数在输入 $y$ 上的值。

\begin{tagblock}{TMs}
\begin{digress}
你可能已经相信，（下点功夫！）可以写一个程序（例如用 Java 或 C++）来完成这件事；现在我们可以诉诸 邱奇--图灵论题，该论题认为任何直观上可计算的事物都能被图灵机计算。

当然，证明 $g(x,y)$ 可计算更直接的方法是，显式地描述一个计算它的图灵机。这尤其能够避免诉诸 邱奇--图灵 论题和直觉。很快，我们将建立起足够的工具，通过一个能被图灵机\emph{模拟}的计算模型来证明 $g(x,y)$ 的可计算性：即递归函数。
\end{digress}
\end{tagblock}

\end{document}